\documentclass[11pt,a4paper]{article}

\usepackage[utf8]{inputenc}
\usepackage[T1]{fontenc}
\usepackage[brazilian]{babel}
\usepackage{lmodern}
\usepackage[margin=2.5cm]{geometry}
\usepackage{amsmath,amssymb}
\usepackage{booktabs}
\usepackage{algorithm}
\usepackage{algpseudocode}
\usepackage{enumitem}
\usepackage{xcolor}
\usepackage{xurl}
\usepackage[hidelinks]{hyperref}

\newcommand{\fd}[2]{\ensuremath{#1 \to #2}}
\newcommand{\cl}[1]{\ensuremath{#1^{+}}}
\newcommand{\cmd}[1]{\texttt{#1}}
\newcommand{\todo}[1]{\textcolor{red}{[TODO: #1]}}
\newcommand{\via}[1]{\xrightarrow{#1}}

\floatname{algorithm}{Algoritmo}
\renewcommand{\listalgorithmname}{Lista de Algoritmos}
\algrenewcommand\algorithmicrequire{\textbf{Entrada:}}
\algrenewcommand\algorithmicensure{\textbf{Saída:}}

\title{fdtool}
\author{Theo S Schult -- GRR20241982}
\date{November 2026}




\begin{document}

\maketitle

\begin{center}
  Repositório: \url{https://github.com/tss6/fdtool/}
\end{center}

%%%%%%%%%%%%%%%%%%%%%%%%%%%%%%%%%%%%%%%%%%%%%%%%%%%%%%%%%%%%%%%%%%%%%%%%%%%%%%%%%%%%%%%%%%%%%%%%%%%%%%
\section{Definições}
%%%%%%%%%%%%%%%%%%%%%%%%%%%%%%%%%%%%%%%%%%%%%%%%%%%%%%%%%%%%%%%%%%%%%%%%%%%%%%%%%%%%%%%%%%%%%%%%%%%%%%

Seja:

\begin{itemize}
    \item $R$: um esquema relacional
    \item $U$: o universo de atributos de $R$
    \item $F$: um conjunto de dependências funcionais sobre $U$
\end{itemize}

\begin{description}[style=nextline]
    \item[Dependência Funcional (DF)]
        Para $X, Y \subseteq U$, $X$ determina $Y$ (\fd{X}{Y}) se cada valor de $X$ é associado com exatamente um valor $Y$.

    \item[DF Trivial]
        \fd{X}{Y} é trivial se $Y \subseteq X$.

    \item[Axiomas de Armstrong]
        Reflexividade (se $Y \subseteq X$ então \fd{X}{Y}), aumento (se \fd{X}{Y}
        então \fd{XZ}{YZ}) e transitividade (se \fd{X}{Y} e \fd{Y}{Z} então
        \fd{X}{Z}).

    \item[Fecho de atributos]
        O fecho \cl{X} de $X$ em relação a $F$ é o maior conjunto $Y \subseteq U$
        tal que \fd{X}{Y} $\in$ \cl{F}.

    \item[Equivalência e cobertura]
        $F$ e $G$ são equivalentes ($F \equiv G$) se $F^{+} = G^{+}$, ou seja, se
        toda DF de $G$ é implicada por $F$ e vice-versa.

    \item[Atributo estranho]
        Em \fd{X}{A} $\in G$, o atributo $B \in X$ é estranho se
        $A \in (X - \{B\})^{+}$ em relação a $G$.

    \item[DF redundante]
        \fd{X}{A} $\in G$ é redundante se $A \in \cl{X}$ em relação a
        $G - \{\fd{X}{A}\}$.

    \item[Cobertura mínima]
        $G$ é uma cobertura mínima de $F$ se $G \equiv F$, todo RHS de $G$ tem um
        único atributo, nenhuma DF de $G$ possui atributo estranho no LHS e
        nenhuma DF de $G$ é redundante. Uma cobertura mínima não é única.

    \item[Superchave]
        $K \subseteq U$ é superchave se $\cl{K} = U$.

    \item[Chave candidata]
        Superchave minimal: $\cl{K} = U$ e, para todo $K' \subsetneq K$,
        $\cl{K'} \neq U$.

    \item[Atributo primo]
        Atributo que pertence a pelo menos uma chave candidata.

    \item[BCNF]
        $R$ está na Forma Normal de Boyce-Codd se, para toda DF não trivial
        \fd{X}{A} de $F^{+}$, $X$ é superchave.

    \item[3FN]
        $R$ está na Terceira Forma Normal se, para toda DF não trivial
        \fd{X}{A} de $F^{+}$, $X$ é superchave \emph{ou} $A$ é primo.
\end{description}

Se uma relação está em BCNF, também está em 3NF. A recíproca não vale.






%%%%%%%%%%%%%%%%%%%%%%%%%%%%%%%%%%%%%%%%%%%%%%%%%%%%%%%%%%%%%%%%%%%%%%%%%%%%%%%%%%%%%%%%%%%%%%%%%%%%%%
\section{Representação}
%%%%%%%%%%%%%%%%%%%%%%%%%%%%%%%%%%%%%%%%%%%%%%%%%%%%%%%%%%%%%%%%%%%%%%%%%%%%%%%%%%%%%%%%%%%%%%%%%%%%%%

\subsection{Bitsets}

Atributos foram restritos a \texttt{A..Z}, escolhi por representar um conjunto de atributos
como um inteiro de 32 bits (\cmd{typedef uint32\_t AttrSet}), onde o bit $i$ determina se o
atributo \texttt{'A'}$+i$ está presente. Dessa forma, as operações de conjunto são operações
binárias (\textit{bitwise}) de custo constante:

\begin{center}
    \begin{tabular}{ll}
        \toprule
        Operação & Implementação \\
        \midrule
        $X \cup Y$ & \cmd{x | y} \\
        $X \cap Y$ & \cmd{x \& y} \\
        $X - Y$ & \cmd{x \& \textasciitilde y} \\
        $X \subseteq Y$ & \cmd{!(x \& \textasciitilde y)} \\
        $|X|$ & \cmd{\_\_builtin\_popcount(x)} (ou \emph{Hamming weight}) \\
        \bottomrule
    \end{tabular}
\end{center}

Uma DF é o par \cmd{struct FD \{ AttrSet lhs; AttrSet rhs; \}} e uma
instância (\cmd{Inst}) é o universo $U$ mais um vetor dinâmico de DFs.

\subsection{Formato \texttt{.fds}}

Desenvolvi um \textit{lexer} e um \textit{parser} para a seguinte
gramática que entendi ser a pretendida:

\begin{verbatim}
file   := decl decl
decl   := 'U' '=' '{' [ attr (sep attr)* ] '}'
        | 'F' '=' '{' [ fd   (sep fd)*   ] '}'
fd     := attr+ '->' attr+
attr   := 'A' .. 'Z'
sep    := ',' | ';'
\end{verbatim}

Qualquer tipo de \textit{whitespace} é ignorado. Ordem não importa.
Erros e avisos são reportados no estilo do compilador GCC.






%%%%%%%%%%%%%%%%%%%%%%%%%%%%%%%%%%%%%%%%%%%%%%%%%%%%%%%%%%%%%%%%%%%%%%%%%%%%%%%%%%%%%%%%%%%%%%%%%%%%%%
\section{Algoritmos}
%%%%%%%%%%%%%%%%%%%%%%%%%%%%%%%%%%%%%%%%%%%%%%%%%%%%%%%%%%%%%%%%%%%%%%%%%%%%%%%%%%%%%%%%%%%%%%%%%%%%%%

Notação para a análise: $n = |U| \le 26$, $m = |F|$ e $m'$ é o número de DFs
após a decomposição do RHS ($m' \le n \cdot m$).



\subsection{Fecho \texorpdfstring{$X^{+}$}{X+}}

\begin{algorithm}[H]
    \caption{Fecho de atributos (\cmd{closure} / \cmd{closure\_skip})}
    \begin{algorithmic}[1]
        \Require conjunto $X$, DFs $F$
        \Ensure \cl{X} em relação a $F$
        \State $mudou \gets 1$
        \While{$mudou$}
        \State $mudou \gets 0$
        \For{$i \gets 0$ \textbf{até} $|F|-1$}
        \If{$\mathit{lhs}_i \subseteq X$ \textbf{e} $\mathit{rhs}_i \not\subseteq X$}
        \State $X \gets X \cup \mathit{rhs}_i$
        \State $mudou \gets 1$
        \EndIf
        \EndFor
        \EndWhile
        \State \Return $X$
    \end{algorithmic}
\end{algorithm}

\paragraph{Explicação.} Cada atributo é adicionado a $X$ a partir de transitividade
com o que já existe em $X$. Se existe uma $\text{DF}_i$ onde podemos aplicar o
axioma, unimos $X$ com o $\mathit{rhs}_i$. O laço \textit{while} termina se nenhuma
DF atualizou $X$, significando que não é possível expandir $X$ além.

\paragraph{Complexidade.} Toda passagem que não é a última acrescenta pelo menos
um atributo, então há no máximo $n+1$ passagens de $m$ testes cada: $O(n \cdot m)$
tempo e $O(1)$ memória extra. Adicionei a variante \cmd{closure\_skip} que ignora
uma DF pelo índice e é usada na remoção de DFs redundantes, evitando copiar o vetor.






\subsection{Cobertura mínima}

\begin{algorithm}[H]
    \caption{Cobertura mínima (\cmd{min\_cover})}
    \begin{algorithmic}[1]
        \Require instância $(U, F)$
        \Ensure cobertura mínima $G$ com RHS unitário
        \State $G \gets \emptyset$
        \ForAll{$\fd{X}{Y} \in F$} \Comment{1. RHS unitário}
        \ForAll{$A \in Y - X$}
        \State $G \gets G \cup \{\fd{X}{A}\}$ \Comment{descarta a parte trivial}
        \EndFor
        \EndFor
        \ForAll{$\fd{X}{A} \in G$} \Comment{2. atributos estranhos}
        \ForAll{$B \in X$ \textbf{com} $|X| > 1$}
        \If{$A \in \cl{(X - \{B\})}$ em relação a $G$}
        \State $X \gets X - \{B\}$
        \EndIf
        \EndFor
        \EndFor
        \ForAll{$\fd{X}{A} \in G$, em ordem} \Comment{3. DFs redundantes}
        \If{$A \in \cl{X}$ em relação a $G - \{\fd{X}{A}\}$}
        \State remove \fd{X}{A} de $G$
        \EndIf
        \EndFor
        \State \Return $G$
    \end{algorithmic}
\end{algorithm}

\paragraph{Explicação.}
\begin{itemize}
    \item O passo 1 aplica a regra da decomposição e descarta os atributos de
        $Y \cap X$, que são triviais. Obviamente, o resultado é equivalente a $F$.
    \item No passo 2, se $A \in \cl{(X-\{B\})}$ então \fd{(X-\{B\})}{A} já é
        implicada por $G$, e trocar \fd{X}{A} por ela mantém a equivalência.
    \item No passo 3, uma DF implicada pelas demais pode ser removida sem
        alterar \cl{G}. A remoção é feita uma DF por vez, recalculando o fecho
        sobre o $G$ já reduzido, o que impede remover duas DFs que só são
        redundantes uma em relação à outra.
\end{itemize}

\paragraph{Complexidade.} O passo 1 custa $O(n \cdot m)$ e gera $m'$ DFs. O
passo 2 faz até $n$ fechos por DF: $O(m' \cdot n \cdot n m') = O(n^2 m'^2)$.
O passo 3 faz um fecho e um deslocamento do vetor (\cmd{memmove}) por DF:
$O(m' \cdot (n m' + m')) = O(n\,m'^2)$. Total: $O(n^2 m'^2)$, ou $O(n^4 m^2)$
em função da entrada. Memória: $O(m')$.






\subsection{Chaves candidatas}

Primeiro classificamos os atributos de $U$ como em grafos direcionados, pela
posição em que aparecem em DFs não triviais de $F$:

\begin{itemize}

    \item \textbf{fonte} ($S$): não aparecem em nenhum RHS (vizinhança de entrada
        vazia). Não podem ser derivados, então pertencem a \emph{toda} chave.

    \item \textbf{sumidouro}: aparecem só em RHS (vizinhança de saída vazia).
        Nunca pertencem a uma chave. Prova: se uma chave $K$ contém um sumidouro
        $Z$, removê-lo não invalida nenhuma DF de ser usada, e $Z$ continua sendo
        derivado, logo $\cl{(K-\{Z\})} = U$ e $K$ não é minima, contradição.

    \item \textbf{meio} ($M$): aparecem em algum LHS e em algum RHS. São os
        únicos que precisam ser testados.
\end{itemize}

Toda chave candidata tem a forma $S \cup T$ com $T \subseteq M$. O algoritmo
enumera os subconjuntos $T$ de $M$ em ordem numérica crescente da máscara de
bits, que parte de $T = \emptyset$ e volta a $\emptyset$ após visitar todas
as $2^{|M|}$ submáscaras.

\begin{algorithm}[H]
    \caption{Chaves candidatas (\cmd{keys})}
    \begin{algorithmic}[1]
        \Require instância $(U, F)$
        \Ensure lista $\mathcal{K}$ de todas as chaves candidatas
        \State $L \gets \bigcup \mathit{lhs}_i$; \ $R \gets \bigcup (\mathit{rhs}_i - \mathit{lhs}_i)$
        \State $S \gets U - R$; \ $M \gets U \cap R \cap L$
        \State $\mathcal{K} \gets \emptyset$; \ $T \gets \emptyset$
        \Repeat
        \State $C \gets S \cup T$
        \If{$\nexists K \in \mathcal{K} : K \subseteq C$ \textbf{e} $\cl{C} = U$}
        \State $\mathcal{K} \gets \mathcal{K} \cup \{C\}$
        \EndIf
        \State $T \gets (T - M) \mathbin{\&} M$ \Comment{próxima submáscara}
        \Until{$T = \emptyset$}
        \State \Return $\mathcal{K}$
    \end{algorithmic}
\end{algorithm}

\paragraph{Explicação (minimalidade).} Se $T' \subsetneq T$, a máscara de $T'$ é
numericamente menor que a de $T$, então todo subconjunto próprio de um
candidato é visitado antes dele. Quando $C$ é testado, todas as chaves contidas
em $C$ já estão em $\mathcal{K}$. Se nenhuma está, nenhum subconjunto próprio
de $C$ é superchave e, se $\cl{C} = U$, $C$ é chave candidata. Se alguma
está, $C$ é superchave não mínima e é descartado.

\paragraph{Complexidade.} São $2^{|M|}$ candidatos. Cada um custa um teste de
minimalidade $O(k)$, sendo $k$ o número de chaves já encontradas, e um fecho
$O(n \cdot m)$. Total: $O(2^{|M|} (n m + k))$, aceitável pois o número de
chaves candidatas pode ser exponencial em $n$. Com $n \le 26$ o pior caso é
$2^{26} \approx 6{,}7 \times 10^7$ candidatos.






\subsection{Verificação de BCNF e 3FN}

\begin{algorithm}[H]
    \caption{Verificação de formas normais (\cmd{nf\_violations})}
    \begin{algorithmic}[1]
        \Require esquema $(U, F)$, chaves candidatas $\mathcal{K}$, modo (BCNF ou 3FN)
        \Ensure lista de DFs que violam a forma normal
        \State $G \gets \{\fd{X}{A} : \fd{X}{Y} \in F,\ A \in Y - X\}$ \Comment{RHS unitário, sem triviais}
        \State $P \gets \bigcup_{K \in \mathcal{K}} K$ \Comment{atributos primos}
        \ForAll{$\fd{X}{A} \in G$}
        \If{$\cl{X} = U$} \textbf{continue} \Comment{$X$ é superchave}
        \EndIf
        \If{modo $=$ 3FN \textbf{e} $A \in P$} \textbf{continue} \Comment{$A$ é primo}
        \EndIf
        \State registra violação \fd{X}{A}
        \EndFor
    \end{algorithmic}
\end{algorithm}

\paragraph{Explicação.} $G$ é $F$ decomposto em RHS unitário, sem as partes
triviais (passo 1 da cobertura mínima). As violações listadas são, portanto,
DFs escritas na entrada. O teste de superchave usa o fecho e o teste de primo
usa a união das chaves candidatas.

Não é preciso examinar todo \cl{F}. Suponha que uma DF não trivial
$\fd{X}{A} \in \cl{F}$ viola a forma normal, ou seja, $X$ não é superchave
(e, na 3FN, $A$ não é primo). Ao calcular \cl{X} com $G$, alguma
$\fd{Y}{A} \in G$ com $Y \subseteq \cl{X}$ adiciona $A$. Se $Y$ fosse
superchave, $U = \cl{Y} \subseteq \cl{X}$ e $X$ seria superchave, contradição.
Logo \fd{Y}{A} também viola a forma normal e é encontrada pelo algoritmo. O
RHS unitário é necessário na 3FN: em \fd{B}{CD} com $C$ primo e $D$ não, só a
parte \fd{B}{D} viola.

\paragraph{Complexidade.} A decomposição custa $O(n m)$ e gera $m'$ DFs. Há
uma chamada a \cmd{keys} (feita uma vez e compartilhada entre os dois modos) e
um fecho por DF de $G$ ($O(m' \cdot n m)$). O custo é dominado por
\cmd{keys}.




\subsection{Resumo das complexidades}

\begin{center}
    \begin{tabular}{lll}
        \toprule
        Operação & Tempo & Memória extra \\
        \midrule
        Leitura e análise do \texttt{.fds} & $O(\text{tamanho do arquivo})$ & $O(\text{tamanho do arquivo})$ \\
        \cmd{closure} & $O(n m)$ & $O(1)$ \\
        \cmd{mincover} & $O(n^2 m'^2) \subseteq O(n^4 m^2)$ & $O(m')$ \\
        \cmd{keys} & $O(2^{|M|}(n m + k))$ & $O(k)$ \\
        \cmd{normalform} & $O(2^{|M|}(n m + k) + n m m')$ & $O(m' + k)$ \\
        \bottomrule
    \end{tabular}
\end{center}






%%%%%%%%%%%%%%%%%%%%%%%%%%%%%%%%%%%%%%%%%%%%%%%%%%%%%%%%%%%%%%%%%%%%%%%%%%%%%%%%%%%%%%%%%%%%%%%%%%%%%%
\section{Estudos de caso}
%%%%%%%%%%%%%%%%%%%%%%%%%%%%%%%%%%%%%%%%%%%%%%%%%%%%%%%%%%%%%%%%%%%%%%%%%%%%%%%%%%%%%%%%%%%%%%%%%%%%%%

Notação: $\{X\} \via{\fd{Y}{Z}} \{XZ\}$ indica que a DF \fd{Y}{Z} foi
aplicada no cálculo de um fecho, e $\cl{X}_{G}$ é o fecho de $X$ em relação
a $G$. As tabelas de formas normais listam as DFs de $F$ com RHS unitário, e
$P$ é o conjunto de atributos primos.

\subsection{Caso 1: exemplo da especificação (violação de BCNF e 3FN)}
\label{sec:caso1}

Arquivo \texttt{exemplos/prof.fds}:
\begin{verbatim}
U = {A,B,C,D,E}
F = { A->BC, C->D, BD->E }
\end{verbatim}

\paragraph{Fecho.}
\begin{align*}
    \cl{A}  &: \{A\} \via{\fd{A}{BC}} \{ABC\} \via{\fd{C}{D}} \{ABCD\}
    \via{\fd{BD}{E}} \{ABCDE\} = U \\
    \cl{BD} &: \{BD\} \via{\fd{BD}{E}} \{BDE\}
\end{align*}

\paragraph{Cobertura mínima.}
\begin{enumerate}
    \item RHS unitário:
        $G = \{\fd{A}{B},\ \fd{A}{C},\ \fd{C}{D},\ \fd{BD}{E}\}$.

    \item Atributos estranhos. Só \fd{BD}{E} tem LHS com mais de um atributo:
        \begin{align*}
            \text{sem } B &: \cl{D}_{G} = \{D\} \not\ni E \\
            \text{sem } D &: \cl{B}_{G} = \{B\} \not\ni E
        \end{align*}
        Nenhum atributo é estranho.

    \item DFs redundantes. Para cada DF, fecho do LHS sem ela:
        \begin{align*}
            \fd{A}{B}  &: \{A\} \via{\fd{A}{C}} \{AC\} \via{\fd{C}{D}} \{ACD\} \not\ni B
            &&\Rightarrow \text{mantém} \\
            \fd{A}{C}  &: \{A\} \via{\fd{A}{B}} \{AB\} \not\ni C
            &&\Rightarrow \text{mantém} \\
            \fd{C}{D}  &: \{C\} \not\ni D
            &&\Rightarrow \text{mantém} \\
            \fd{BD}{E} &: \{BD\} \not\ni E
            &&\Rightarrow \text{mantém}
        \end{align*}
\end{enumerate}
$G$ não muda:
\begin{verbatim}
$ ./fdtool mincover --fds exemplos/prof.fds
A->B
A->C
C->D
BD->E
\end{verbatim}

\paragraph{Chaves.}
\begin{align*}
    L &= A \cup C \cup BD = \{ABCD\} \\
    R &= BC \cup D \cup E = \{BCDE\} \\
    S &= U - R = \{A\}               && \text{(fonte)} \\
    M &= U \cap R \cap L = \{BCD\}   && \text{(meio)} \\
    R - L &= \{E\}                   && \text{(sumidouro)}
\end{align*}
O primeiro candidato é $S \cup \emptyset = \{A\}$, e $\cl{A} = U$ (fecho
acima), então $\{A\}$ é chave. Os outros $2^3 - 1$ candidatos contêm $\{A\}$
e são descartados sem calcular fecho. $\mathcal{K} = \{A\}$ e $P = \{A\}$.

\paragraph{Formas normais.}
\begin{center}
    \begin{tabular}{lllccc}
        \toprule
        \fd{X}{A} & \cl{X} & $X$ superchave? & $A \in P$? & BCNF & 3FN \\
        \midrule
        \fd{A}{B}  & $\{ABCDE\} = U$ & sim & não & $\checkmark$ & $\checkmark$ \\
        \fd{A}{C}  & $\{ABCDE\} = U$ & sim & não & $\checkmark$ & $\checkmark$ \\
        \fd{C}{D}  & $\{CD\}$        & não & não & $\times$     & $\times$     \\
        \fd{BD}{E} & $\{BDE\}$       & não & não & $\times$     & $\times$     \\
        \bottomrule
    \end{tabular}
\end{center}
\begin{verbatim}
$ ./fdtool normalform --fds exemplos/prof.fds
BCNF: VIOLATIONS
3NF: VIOLATIONS
VIOLATION BCNF: C->D (C not superkey)
VIOLATION BCNF: BD->E (BD not superkey)
VIOLATION 3NF: C->D (D not prime, C not superkey)
VIOLATION 3NF: BD->E (E not prime, BD not superkey)
\end{verbatim}
\fd{C}{D} e \fd{BD}{E} são dependências transitivas da chave:
$A \to C \to D$ e $A \to BD \to E$.





\subsection{Caso 2: em 3FN, mas não em BCNF}
\label{sec:caso2}

Arquivo \texttt{exemplos/multkeys.fds}:
\begin{verbatim}
U = {A,B,C}
F = { AB->C, C->B }
\end{verbatim}

\paragraph{Cobertura mínima.}
\begin{enumerate}
    \item RHS unitário: $G = F = \{\fd{AB}{C},\ \fd{C}{B}\}$.

    \item Atributos estranhos em \fd{AB}{C}:
        \begin{align*}
            \text{sem } A &: \cl{B}_{G} = \{B\} \not\ni C \\
            \text{sem } B &: \cl{A}_{G} = \{A\} \not\ni C
        \end{align*}

    \item DFs redundantes:
        \begin{align*}
            \fd{AB}{C} &: \{AB\} \not\ni C &&\Rightarrow \text{mantém} \\
            \fd{C}{B}  &: \{C\} \not\ni B  &&\Rightarrow \text{mantém}
        \end{align*}
\end{enumerate}
$G$ não muda. Saída: \texttt{AB->C}, \texttt{C->B}.



\paragraph{Chaves.}
\begin{align*}
    L &= AB \cup C = \{ABC\} \\
    R &= C \cup B = \{BC\} \\
    S &= U - R = \{A\}             && \text{(fonte)} \\
    M &= U \cap R \cap L = \{BC\}  && \text{(meio)} \\
    R - L &= \emptyset             && \text{(sumidouro)}
\end{align*}
Enumeração dos subconjuntos $T \subseteq M$ em ordem crescente de máscara:
\begin{center}
    \begin{tabular}{llll}
        \toprule
        $T$ & $S \cup T$ & Fecho & Resultado \\
        \midrule
        $\emptyset$ & $\{A\}$   & $\{A\}$
        & não é superchave \\
        $\{B\}$     & $\{AB\}$  & $\{AB\} \via{\fd{AB}{C}} \{ABC\} = U$
        & chave \\
        $\{C\}$     & $\{AC\}$  & $\{AC\} \via{\fd{C}{B}} \{ABC\} = U$
        & chave \\
        $\{BC\}$    & $\{ABC\}$ & --
        & contém $\{AB\}$, descartado \\
        \bottomrule
    \end{tabular}
\end{center}
$\mathcal{K} = \{AB,\ AC\}$ e $P = AB \cup AC = \{ABC\} = U$.



\paragraph{Formas normais.}
\begin{center}
    \begin{tabular}{lllccc}
        \toprule
        \fd{X}{A} & \cl{X} & $X$ superchave? & $A \in P$? & BCNF & 3FN \\
        \midrule
        \fd{AB}{C} & $\{ABC\} = U$ & sim & sim & $\checkmark$ & $\checkmark$ \\
        \fd{C}{B}  & $\{BC\}$      & não & sim & $\times$     & $\checkmark$ \\
        \bottomrule
    \end{tabular}
\end{center}
\begin{verbatim}
$ ./fdtool normalform --fds exemplos/multkeys.fds
BCNF: VIOLATIONS
3NF: OK
VIOLATION BCNF: C->B (C not superkey)
\end{verbatim}
\fd{C}{B} viola a BCNF porque $C$ não é superchave, mas satisfaz a 3FN porque
$B$ é primo.






\subsection{Caso 3: redundância e atributo estranho}
\label{sec:caso3}

Arquivo \texttt{exemplos/extra.fds}:
\begin{verbatim}
U = {A,B,C,D}
F = { A->B, B->C, A->C, AB->D }
\end{verbatim}

\paragraph{Fecho.}
\begin{align*}
    \cl{A} &: \{A\} \via{\fd{A}{B}} \{AB\} \via{\fd{B}{C}} \{ABC\}
    \via{\fd{AB}{D}} \{ABCD\} = U \\
    \cl{B} &: \{B\} \via{\fd{B}{C}} \{BC\}
\end{align*}

\paragraph{Cobertura mínima.}
\begin{enumerate}
    \item RHS unitário: todas as DFs já têm RHS unitário e nenhuma é trivial,
        então $G = F$.

    \item Atributos estranhos. Só \fd{AB}{D} tem LHS com mais de um atributo:
        \begin{align*}
            \text{sem } A &: \cl{B}_{G} = \{BC\} \not\ni D \\
            \text{sem } B &: \cl{A}_{G} = \{ABCD\} \ni D
            &&\Rightarrow B \text{ é estranho, } \fd{AB}{D} \text{ vira } \fd{A}{D}
        \end{align*}
        Agora $G = \{\fd{A}{B},\ \fd{B}{C},\ \fd{A}{C},\ \fd{A}{D}\}$. O
        atributo $B$ é estranho porque $A$ já determina $B$.

    \item DFs redundantes, em ordem. O fecho é calculado sobre o $G$ já
        reduzido, sem a DF testada:
        \begin{align*}
            \fd{A}{B} &: \{A\} \via{\fd{A}{C}} \{AC\} \via{\fd{A}{D}} \{ACD\} \not\ni B
            &&\Rightarrow \text{mantém} \\
            \fd{B}{C} &: \{B\} \not\ni C
            &&\Rightarrow \text{mantém} \\
            \fd{A}{C} &: \{A\} \via{\fd{A}{B}} \{AB\} \via{\fd{B}{C}} \{ABC\} \ni C
            &&\Rightarrow \text{remove} \\
            \fd{A}{D} &: \{A\} \via{\fd{A}{B}} \{AB\} \via{\fd{B}{C}} \{ABC\} \not\ni D
            &&\Rightarrow \text{mantém}
        \end{align*}
\end{enumerate}
\begin{verbatim}
$ ./fdtool mincover --fds exemplos/extra.fds
A->B
B->C
A->D
\end{verbatim}

\paragraph{Chaves.}
\begin{align*}
    L &= A \cup B \cup A \cup AB = \{AB\} \\
    R &= B \cup C \cup C \cup D = \{BCD\} \\
    S &= U - R = \{A\}            && \text{(fonte)} \\
    M &= U \cap R \cap L = \{B\}  && \text{(meio)} \\
    R - L &= \{CD\}               && \text{(sumidouros)}
\end{align*}
O primeiro candidato é $S \cup \emptyset = \{A\}$, e $\cl{A} = U$ (fecho
acima), então $\{A\}$ é chave. O único outro candidato, $\{AB\}$, contém
$\{A\}$ e é descartado. $\mathcal{K} = \{A\}$ e $P = \{A\}$.

\paragraph{Formas normais.}
\begin{center}
    \begin{tabular}{lllccc}
        \toprule
        \fd{X}{A} & \cl{X} & $X$ superchave? & $A \in P$? & BCNF & 3FN \\
        \midrule
        \fd{A}{B}  & $\{ABCD\} = U$ & sim & não & $\checkmark$ & $\checkmark$ \\
        \fd{B}{C}  & $\{BC\}$       & não & não & $\times$     & $\times$     \\
        \fd{A}{C}  & $\{ABCD\} = U$ & sim & não & $\checkmark$ & $\checkmark$ \\
        \fd{AB}{D} & $\{ABCD\} = U$ & sim & não & $\checkmark$ & $\checkmark$ \\
        \bottomrule
    \end{tabular}
\end{center}
\begin{verbatim}
$ ./fdtool normalform --fds exemplos/extra.fds
BCNF: VIOLATIONS
3NF: VIOLATIONS
VIOLATION BCNF: B->C (B not superkey)
VIOLATION 3NF: B->C (C not prime, B not superkey)
\end{verbatim}
A tabela inclui \fd{A}{C} e \fd{AB}{D}, embora a cobertura mínima remova a
primeira e reescreva a segunda, pois a verificação usa as DFs da entrada.
Ambas têm LHS superchave. \fd{B}{C} é uma dependência transitiva da chave:
$A \to B \to C$.





%%%%%%%%%%%%%%%%%%%%%%%%%%%%%%%%%%%%%%%%%%%%%%%%%%%%%%%%%%%%%%%%%%%%%%%%%%%%%%%%%%%%%%%%%%%%%%%%%%%%%%
\section{Limitações}
%%%%%%%%%%%%%%%%%%%%%%%%%%%%%%%%%%%%%%%%%%%%%%%%%%%%%%%%%%%%%%%%%%%%%%%%%%%%%%%%%%%%%%%%%%%%%%%%%%%%%%

\begin{itemize}
    \item \textbf{Universo limitado a 26 atributos} de uma letra
        (\texttt{A..Z}, sem minúsculas), o que foi o motivo principal de
        bitsets terem sido escolhidos em vez de grafos.

    \item \textbf{Custo exponencial de \cmd{keys}} no número de atributos do
        ``meio'' ($2^{|M|}$, até $2^{26}$). As entradas tendem a ser pequenas
        devido a limitação do Universo, mas instâncias muito grandes podem
        demorar um pouco.

    \item \textbf{Cobertura mínima dependente da ordem} da entrada.
        Os arquivos \texttt{order1.fds} e \texttt{order2.fds} diferem apenas na
        ordem de \fd{A}{C} e \fd{B}{C} e produzem coberturas diferentes
        (\fd{B}{C} contra \fd{A}{C}), ambas válidas. Ideal seria identificar
        e potencialmente imprimir todas as coberturas que a instância permite.

    \item \textbf{Deduplicação apenas na saída.} Por exemplo, na sequinte
        dependência: $U = \{A,B,C\}$ e $F = \{\fd{A}{B},\ \fd{B}{C},\ \fd{B}{C}\}$,
        tudo funciona sem \textit{warnings}, porém a saída imprime apenas uma vez
        cada FD única, como explicado no e-mail.

    \item \textbf{Sem decomposição}: a ferramenta diagnostica as formas normais,
        mas não propõe decomposições em BCNF ou 3FN.

\end{itemize}





%%%%%%%%%%%%%%%%%%%%%%%%%%%%%%%%%%%%%%%%%%%%%%%%%%%%%%%%%%%%%%%%%%%%%%%%%%%%%%%%%%%%%%%%%%%%%%%%%%%%%%
\section{Principais decisões}
%%%%%%%%%%%%%%%%%%%%%%%%%%%%%%%%%%%%%%%%%%%%%%%%%%%%%%%%%%%%%%%%%%%%%%%%%%%%%%%%%%%%%%%%%%%%%%%%%%%%%%

\begin{itemize}
    \item \textbf{Bitset para conjuntos de atributos.} Foi considerada uma
        representação em grafo, mas com no máximo 26 atributos operações de
        conjuntos ficam simples e rápidas.

    \item \textbf{Analisador léxico e sintático} com base na gramática que
        esperada e diagnósticos no estilo do GCC. O subcomando extra,
        \cmd{check}, só valida o arquivo.

    \item \textbf{Enumeração crescente de submáscaras} em \cmd{keys},o que
        garante minimalidade sem um passo extra de filtragem e reduz o espaço
        de busca.

    \item \textbf{Interface.} Opções aceitas como \texttt{----fds arq} ou
        \texttt{----fds=arq}, em qualquer ordem; \texttt{----X} aceita \texttt{AB},
        \texttt{A,B} ou \texttt{A B} e é validado contra $U$. Códigos de saída:
        0 (ok), 1 (erro de leitura) e 2 (uso incorreto).

    \item \textbf{Saída ordenada.} Conjuntos são impressos percorrendo os bits
        de \texttt{A} a \texttt{Z}, ordenando automaticamente.
\end{itemize}





%%%%%%%%%%%%%%%%%%%%%%%%%%%%%%%%%%%%%%%%%%%%%%%%%%%%%%%%%%%%%%%%%%%%%%%%%%%%%%%%%%%%%%%%%%%%%%%%%%%%%%
\section{Uso de IA e referências}
%%%%%%%%%%%%%%%%%%%%%%%%%%%%%%%%%%%%%%%%%%%%%%%%%%%%%%%%%%%%%%%%%%%%%%%%%%%%%%%%%%%%%%%%%%%%%%%%%%%%%%

Inteligência artificial foi utilizada para pesquisar tópicos e esclarecer conceitos, não
para escrever código-fonte. A versão base deste relatório foi gerada com auxílio de IA
a partir do código e dos exemplos do repositório, e depois revisada pelo autor.

{\small
\begin{itemize}[nosep]
    \item \textbf{Disciplina}
    \begin{itemize}[nosep]
        \item Slides da disciplina de Banco de Dados.
    \end{itemize}
    \item \textbf{Dependências funcionais e cobertura mínima}
    \begin{itemize}[nosep]
        \item \url{https://www.geeksforgeeks.org/dbms/how-to-find-minimal-set/}
        \item \url{https://stackoverflow.com/questions/10284004/}
    \end{itemize}
    \item \textbf{Enumeração de submáscaras}
    \begin{itemize}[nosep]
        \item \url{https://www.inf.ufpr.br/andre/Disciplinas/CI1031-2026-2/03-busca-completa-matematica.pdf}
    \end{itemize}
    \item \textbf{Contagem de bits}
    \begin{itemize}[nosep]
        \item \url{https://stackoverflow.com/questions/109023/}
        \item \url{https://www.felixcloutier.com/x86/popcnt}
        \item \url{https://gcc.gnu.org/onlinedocs/gcc/Bit-Operation-Builtins.html#index-_005f_005fbuiltin_005fpopcount}
        \item \url{https://clang.llvm.org/docs/LanguageExtensions.html#builtin-popcountg}
        \item \url{https://learn.microsoft.com/en-us/cpp/intrinsics/popcnt16-popcnt-popcnt64?view=msvc-170}
        \item \url{https://en.wikipedia.org/wiki/Hamming_weight}
    \end{itemize}
    \item \textbf{Vetor dinâmico}
    \begin{itemize}[nosep]
        \item \url{https://www.youtube.com/watch?v=gtk3RZHwJUA}
        \item \url{https://stackoverflow.com/questions/1100311/}
    \end{itemize}
    \item \textbf{Analisador léxico e sintático}
    \begin{itemize}[nosep]
        \item \url{https://github.com/lotabout/write-a-C-interpreter/blob/master/tutorial/en/3-Lexer.md}
        \item \url{https://stackoverflow.com/questions/17975913/}
        \item \url{https://github.com/rexim/bnfuzzer}
    \end{itemize}
    \item \textbf{Linguagem C}
    \begin{itemize}[nosep]
        \item \url{https://stackoverflow.com/questions/7762731/}
        \item \url{https://sourceware.org/glibc/manual/latest/html_node/Feature-Test-Macros.html}
        \item \url{https://gcc.gnu.org/onlinedocs/gcc/Common-Attributes.html#index-format}
        \item \url{https://stackoverflow.com/a/35281548}
    \end{itemize}
    \item \textbf{Cores no terminal}
    \begin{itemize}[nosep]
        \item \url{https://gist.github.com/JBlond/2fea43a3049b38287e5e9cefc87b2124}
    \end{itemize}
\end{itemize}
}

\end{document}
